\chapter{催化日历：从H1落地到年末验证}

\section{催化必须在时间窗内发生}

本题的截止日是2026年12月31日，因此长期产能、2028年矿山投产或2030年订单能见度只能解释资产质量，不能替代年末催化。真正有用的催化集中在H1正式报告、三季报、产品价格与出货、经营现金流以及下一年度预期形成。

\begin{exhibitbox}[Exhibit 9：2026年下半年验证时间轴]
\centering
\begin{tikzpicture}[x=1cm,y=1cm,font=\small]
\draw[very thick,navy] (0,0)--(14,0);
\foreach \x/\lab in {0/7月下旬,3/8月,6/9月,9/10月,12/11月,14/12月末}{
  \draw[navy] (\x,0.18)--(\x,-0.18);
  \node[below=0.25cm] at (\x,0) {\lab};
}
\node[above=0.45cm,align=center,text width=2.8cm] at (2,0) {H1正式报告\\预告质量复核};
\node[above=1.35cm,align=center,text width=3cm] at (6,0) {锂盐/隔膜价格、出货\\新线与库存证据};
\node[above=0.45cm,align=center,text width=3cm] at (9,0) {三季报\\利润与OCF验证};
\node[above=1.35cm,align=center,text width=3cm] at (12.5,0) {2027预期形成\\年末估值选择};
\end{tikzpicture}
\sourcenote{AStock研究框架；具体定期报告日期以交易所安排为准。}
\end{exhibitbox}

\section{雅化集团：最短验证路径}

H1正式报告首先需要解释利润来源：销量、实现价、自供矿、库存与套保各自贡献多少。三季报需要看到锂盐出货接近3万吨、剔除库存与套保后的吨利至少约2万元、经营现金流转正或大幅改善。到年末，市场还需要相信全年约12万吨出货、约30亿元利润与民爆5.5--6亿元稳定利润能够持续。

如果利润兑现但现金仍差，基准目标可以成立，翻倍概率不应上调；如果利润、现金、自供比例和合同履行同时改善，12.77倍年末PE才有合理基础。

\section{恩捷股份：验证2027而非只验证H1}

H1正式报告只能确认扭亏。决定103元目标的是2027年盈利跃升，因此Q3必须出现客户涨价、出货约40亿平方米、单平利润约0.15元、新线良率与经营现金流等公司级证据。若这些参数仍只存在于券商模型，市场可能继续把2027E EPS按低倍数折现。

\section{中国船舶、盛新锂能与江波龙}

中国船舶需要高价船交付、毛利率约16\%以上和正经营现金流；盛新锂能需要约3万吨季度出货、约2万元综合吨利、印尼利用率与正现金流；江波龙需要单季现金转正、库存周转改善、毛利稳定和2027年增长重新加速。

\begin{exhibitbox}[Exhibit 10：升级与删除条件]
\small
\begin{tabularx}{\exhibitboxwidth}{L{2.0cm} X L{2.4cm} X}
\toprule
标的 & 升级证据 & 阈值性质 & 删除/降级证据 \\
\midrule
雅化集团 & H2利润桥、OCF转正、自供与民爆稳定 & AStock阈值；部分输入源于券商 & 现金持续为负、库存收益主导、锂价反转 \\
恩捷股份 & 客户涨价、单平利润、利用率/良率与OCF & 券商估计转为公司证据 & 单平低于约0.12元、涨价延迟、扩产吞噬现金 \\
中国船舶 & 高价交付、毛利、OCF及显式外部目标 & AStock阈值；经营输入源于券商 & 交付延迟、毛利回落、整合未兑现 \\
盛新锂能 & 出货、吨利、自有矿、印尼与OCF共同验证 & 券商估计/AStock阈值 & 锂价回落、资源贡献不足、短债/现金恶化 \\
江波龙 & 现金转正、库存改善、2027利润增长两位数 & AStock阈值；预测源于券商 & 毛利回落、库存继续上升、现金仍显著为负 \\
天华新能 & 资源自供与OCF提升且价格回到安全区间 & AStock阈值 & 牛市价值仍低于翻倍门槛 \\
\bottomrule
\end{tabularx}
\qualitynote{券商派生阈值不是公司指引或已实现事实；只有定期报告、公告或可审计经营数据能关闭验证缺口。}
\end{exhibitbox}

\clearpage
